%% Chapter 18: Kernelization
%% Status: skeleton.  Only the learning goals are fixed; they come from the
%% exam pool.  The sectioning is deliberately not invented in advance --
%% it will follow the source material once that has been read.

\chapter{Kernelization}
\label{ch:kernelization}

Kernelization is preprocessing with a guarantee. We already saw a quadratic
kernel for \lang{Vertex Cover}; here we push it to linear using linear
programming, develop the general tools, and then show that for some problems no
polynomial kernel exists at all.

\begin{goals}
  \poolgoal{define \lang{Integer Linear Programming} and show how \lang{Vertex Cover} can be encoded as one}
  \goalitem{prove that the LP relaxation of \lang{Vertex Cover} is half-integral}
  \poolgoal{show that \lang{Vertex Cover} has a linear sized kernel}
  \goalitem{use crown decompositions and the sunflower lemma to obtain kernels}
  \goalitem{explain how composition arguments rule out polynomial kernels}
\end{goals}

\todo{This chapter is not written yet. Write it against
\texttt{material/slides/aa-fpt.pdf}.}
